\documentclass[11pt,a4paper]{article}
\usepackage{gtu-el-yazisi}
\begin{document}
\notbilgi{Analog Haberleşme}{Süperheterodin alıcı, açı modülasyonu, NBFM}{Ders11.pdf}

\sayfa{1}
\section{AM Radio Broadcasting (GM Radio yayını)}
\hfill (\okunamadi{BT-448})

Genlik Modülasyonu (GM) Radyo Alıcısı\\
(süper heterodin alıcı -- Superheterodyne receiver)

Intermediate frequency (IF) amplifier\\
$f_{LO}$: Local oscillator frequency
\[ f_{LO}=f_c+f_{IF} \qquad\qquad f_c'=f_{LO}+f_{IF} \]

\sekilnotu[1]{Bu sayfa basılı materyal içeriyor (Figure 3.34: Frequency-response characteristics of both IF and RF amplifiers; (a) $f_{IF}$ merkezli $B_c$ genişlikli IF bandı, (b) $f_c=f_{LO}-f_{IF}$ merkezli $B_{RF}$ genişlikli RF bandı, (c) $f_c$ ve $f_c'=f_{LO}+f_{IF}$ konumlarında iki bant, aralarında $f_{LO}$ ve $2f_{IF}$); dijitalleştirilmedi. El yazısı eklemeler aşağıda aktarılmıştır; (c) şeklinin altına ayrıca ``910 kHz'' yazılmış.}

\[ B_c < B_{RF} < 2f_{IF} \quad\text{(altlarında sırasıyla: } 10\ \mathrm{kHz},\ 910\ \mathrm{kHz}\text{)} \]
\[ 455\ \mathrm{kHz} = f_{IF} \quad\text{(intermediate frequency)} \]
\[ \text{e.g.}\quad f_c=1000,\quad f_{LO}=1455,\quad f_c'=1910,\quad f_c'-f_c=910\ \mathrm{kHz} \]

GM frekans yelpazesi: 535--1605 kHz\\
GM taşıyıcı frekans yelpazesi: 540--1600 kHz
\[ B=\mathrm{BG}\{m(t)\}=5\ \mathrm{kHz} \;\Rightarrow\; \mathrm{BG}\{y_{\text{DSB-LC-AM}}(t)\}=2B=10\ \mathrm{kHz} \]
1. istasyon:
\[ y_1(t)=A_c\,[1+\mu\, m_{n1}(t)]\cos(2\pi f_c t) \]
2. istasyon:
\[ y_2(t)=A_c\,[1+\mu\, m_{n2}(t)]\cos(2\pi f_c' t), \qquad f_c' :\ \text{Image station} \]
RF amplifier sadece $y_1$'i geçirir.

\elcizimi[0.4]{sekiller/Ders11-s1-1.png}[1]

\sayfa{2}
\sekilnotu[2]{Bu sayfa basılı materyal içeriyor (Figure 3.33: A superheterodyne receiver); dijitalleştirilmedi. Basılı blok diyagram: anten $\to$ RF amplifier $\to$ Mixer $\to$ IF amplifier $\to$ Detector $\to$ Audio frequency amplifier $\to$ Speaker; Local oscillator Mixer'e bağlı (RF amplifier'a kesikli çizgiyle de bağlı); Detector'dan IF amplifier'a Automatic volume control (AVC) geri beslemesi. El yazısı eklemeler: RF amplifier ve IF amplifier üzerinde (BGF); RF amplifier altında $f_c=f_{LO}-f_{IF}$; IF amplifier altında $f_{IF}=455$ kHz, $B=10$ kHz; RF amplifier çıkışı / Mixer girişi: $A_c[1+\mu m_n(t)]\cos(2\pi f_c t)$; Mixer çıkışı / IF girişi: $A_c[1+\mu m_{n1}(t)]\cos(2\pi f_{IF} t)+2f_{IF}$ frekanslı bileşen (``$2f_{IF}$ frekanslı bileşen'' kısmı turuncu çizgiyle iptal edilmiş); IF çıkışı / Detector girişi: $A_c[1+\mu m_{n1}(t)]\cos(2\pi f_{IF} t)$; hoparlör girişi: $K\,m_1(t)$; Local oscillator altında $f_{LO}=f_c+f_{IF}$ ve ``Tuning range (Ayarlama Aralığı): 995--2055 kHz''.}

RF kuvvetlendiricisinin çıkışında
\begin{align*}
y_1(t) &= A_c\,[1+\mu\, m_{n1}(t)]\cos(2\pi f_c t)\\
y_2(t) &= A_c\,[1+\mu\, m_{n2}(t)]\cos(2\pi f_c' t) \quad\text{olsun}
\end{align*}
\[ f_c=\overset{\text{local osilator}}{f_{LO}}-f_{IF} \;\to\; \text{Intermediate frequency} \]
\[ f_c'=f_{LO}+f_{IF} \qquad\text{ve}\qquad f_{LO}=f_c+f_{IF} \]
mixer çıkışı:
\begin{align*}
z_1(t) &= A_c\,[1+\mu\, m_{n1}(t)]\cos(2\pi f_{IF}(t)) + 2f_c\ \text{frekanslı bileşen}\\
z_2(t) &= A_c\,[1+\mu\, m_{n2}(t)]\cos(2\pi f_{IF}(t)) + 2f_c\ \text{frekanslı bileşen}
\end{align*}
[$z_2$ satırının sonuna okla işaret edilmiş turuncu not: \okunamadi{$f_c'$'den gelen}]

IF amplifier'ın girişine image station ($f_c'$) gelmemeli yoksa çıkışta işaretler birbirine karışıyor. Bunun için RF amplifier'ın görevi dinlemek istediğimiz istasyonun image istasyonunu geçirmemektir.

IF amplifier'ın görevi de dinlemek istemediğimiz diğer frekansları geçirmemektir. (IF girişine dinlemek istemediğimiz frekanslar da geliyor.)

Çok yüksek frekansta olduğu için ilk başta direkt 10 kHz'le BGF yapamıyoruz. Bu kadar yüksek frekanslarda 10 kHz seçiciliği olan bir BGF tasarlanamıyor.

\sayfa{3}
\section{Açı (Angle) Modülasyonu}
\begin{itemize}
\item Lineer değil
\item Yüksek BG $\to$ Yüksek gürültü bağışıklığı
\item Sabit Zarflı ($A_c$'den dolayı) $\to$ Modülasyon tekniklerinin sabit zarflı olması tercih edilir. Çünkü sabit zarflı olmazsa devre doymaya girebilir. Bu da sinyalde kesintiye neden olur.
\end{itemize}
\[ y(t)=A_c\cos(\theta(t)) \]
\begin{align*}
\theta(t)&=2\pi f_c t+\underset{\text{sabit}}{\theta_0} \;\to\; \underset{\substack{\text{Ani frekans}\\ \text{(Instaneous frequency)}}}{f_i(t)}=\underbrace{\frac{1}{2\pi}\frac{d\theta(t)}{dt}}_{f_c}\\
&\;\to\; \theta(t)=2\pi\int_{-\infty}^{t} f_i(\alpha)\,d\alpha
\end{align*}
ifadeyi genelleştirelim
\begin{align*}
\theta(t)&=2\pi f_c t+\underset{\substack{\text{mesaj işaretini}\\ \text{faza koyuyoruz}}}{\overset{\text{Faz}}{\phi(t)}} \;\to\; f_i(t)=f_c+\frac{1}{2\pi}\frac{d\phi(t)}{dt}\\
&\Rightarrow\ \text{Ani frekans zamanın fonksiyonu oldu}
\end{align*}
\begin{align*}
\underset{\substack{\text{modüle ettiğimiz}\\ \text{işaretin ani frekansı}}}{f_i(t)}-\underset{\substack{\text{taşıyıcı}\\ \text{frekansı}}}{f_c}&=\frac{1}{2\pi}\frac{d\phi(t)}{dt}\\
&\Rightarrow\ \text{Ani frekansın taşıyıcıdan sapma miktarı}
\end{align*}

\subsection{1) Phase Mod. (PM) / Faz Modülasyonu (FZM)}
\[ \phi(t)=k_p\, m(t) \qquad\qquad \theta(t)=2\pi f_c t+k_p\, m(t) \]
($k_p$: faz sapma sabiti (Phase deviation constant))
\[ y_{\mathrm{FZM}}(t)=A_c\cos\Big(\overbrace{2\pi f_c t+k_p\, m(t)}^{\theta(t)}\Big) \]
\[ f_i(t)=f_c+\frac{1}{2\pi}\,k_p\,\frac{dm(t)}{dt} \]

\sayfa{4}
\subsection{2) Frequency Mod. (FM) / Frekans Modülasyonu (FRM)}
\[ f_i(t)=f_c+k_f\, m(t) \;\Rightarrow\; \text{ani frekans, mesaj işaretiyle lineer olarak değişiyor.} \]
($k_f$: frekans sapma sabiti (frequency deviation constant))
\begin{align*}
\Rightarrow\ \theta(t)&=2\pi f_c t+2\pi k_f\int_{-\infty}^{t} m(\tau)\,d\tau\\
y_{\mathrm{FRM}}(t)&=A_c\cos\Big(2\pi f_c t+2\pi k_f\int_{-\infty}^{t} m(\tau)\,d\tau\Big)\\
y_{\text{Açı}}(t)&=A_c\cos\big(2\pi f_c t+\phi(t)\big)
\end{align*}
\begin{align*}
&\hookrightarrow\ {=}\,k_p\, m(t) \quad \text{FZM}\\
&\hookrightarrow\ {=}\,2\pi k_f\int_{-\infty}^{t} m(\tau)\,d\tau \quad \text{FRM}
\end{align*}

\begin{center}
\begin{tikzpicture}[blok/.style={draw=murekkep,thick,minimum height=10mm,align=center,text=murekkep}]
\node[left] (in1) at (0,0) {$m(t)$};
\node[blok,minimum width=10mm] (int) at (1.8,0) {$\displaystyle\int$};
\node[blok] (pm) at (5,0) {PM Modulator\\ (FZM)};
\draw[->] (in1.east) -- (int.west);
\draw[->] (int.east) -- (pm.west);
\draw[->] (pm.east) -- ++(1.8,0) node[midway,above] {FRM} node[midway,below] {$y_{\mathrm{FRM}}(t)$};
\end{tikzpicture}

\bigskip
\begin{tikzpicture}[blok/.style={draw=murekkep,thick,minimum height=10mm,align=center,text=murekkep}]
\node[left] (in2) at (0,0) {$m(t)$};
\node[blok,minimum width=10mm] (dif) at (1.8,0) {$\dfrac{d}{dt}$};
\node[blok] (fm) at (5,0) {FM Modulator\\ (FRM)};
\draw[->] (in2.east) -- (dif.west);
\draw[->] (dif.east) -- (fm.west) node[midway,below] {$\dfrac{dm(t)}{dt}$};
\draw[->] (fm.east) -- ++(1.8,0) node[midway,above] {FZM} node[midway,below] {$y_{\mathrm{FZM}}(t)$};
\end{tikzpicture}
\end{center}

\underline{En genel halde}\\
Mesaj işaretinin genliği ne olursa olsun bu modülasyon tekniğinin gücü: $P_y=\dfrac{A_c^2}{2}$

\subsection{Genel üstel (Açı) Mod.}
\begin{align*}
y_{\text{ÜM}}(t)&=A_c\cos\big(2\pi f_c t+\phi(t)\big)\\
\phi(t)&=\int_{-\infty}^{t} m(\tau)\,h(t-\tau)\,d\tau
\end{align*}
Eğer $h(t)=k_p\,\delta(t)$ $\to$ FZM elde ederiz\\
Eğer $h(t)=2\pi k_f\,u(t)$ $\to$ FRM elde ederiz.

Optimum performans ikisinin arasında sağlanır. Ancak bunu gerçekleştirmeye kalktığımızda hem verici hem alıcı karmaşıklaşıyor. Bu da maliyet vs. arttırıyor. Biz Radyo yayını yaptığımızdan alıcının basit olmasını istiyoruz.

\medskip
\noindent\underline{Analog}\\
5 Ocak Cuma 11.00\\
$\hookrightarrow$ Bonus Quiz\\
Buradan alınan puan ara sınava eklenecek.

\sayfa{5}
\sekilnotu[5]{Bu sayfa basılı materyal içeriyor (Figure 4.2: Frequency and phase modulations of square and sawtooth waves); dijitalleştirilmedi. Şekil: sol üstte kare dalga $m(t)$ ($0$--$2$ arası $+1$, $2$--$4$ arası $-1$), sağ üstte $\int_{-\infty}^{t} m(\tau)\,d\tau$ (üçgen dalga: $t=2$'de $2$'ye çıkıp $t=4$'te $0$'a iniyor); altta solda FM signal ve PM signal, sağda ``Equivalent'' okuyla eşlenmiş FM signal ve PM signal grafikleri. El yazısı (kırmızı) eklemeler: $m(t)$ grafiğinde $t=2$ noktasına okla ``işaret(?) değişiyor''; sol FM grafiğinin solunda ``YF: Yüksek frekans'', ``DF: Düşük frekans'', altında $t\approx1$ için ``YF $f_c+k_f$'', $t\approx3$ için ``DF $f_c-k_f$''; sağ üst FM grafiğinin altında $0$--$1$ arası ``DF'', $1$--$3$ arası ``YF'' (parantezle), $3$--$4$ arası ``DF''; sol PM grafiğinin solunda ``frekans hiç değişmiyor'', ``faz mesajın(?) değiştiği yerde değişiyor'' ve $t=2$'de küçük kırmızı işaret; sağ alt PM grafiğinin altında $0$--$2$ arası ``$f_i(t)=f_c+k_f$'', $2$--$4$ arası ``$f_c-k_f$'' ve şeklin sağ altındaki formüllere kırmızı okla işaret ediliyor.}
\[ f_i(t)=f_c+\overbrace{\frac{k_p}{2\pi}}^{k_f}\,\frac{dm(t)}{dt} \]
\[ f_i(t)=f_c+k_f\,\frac{dm(t)}{dt} \]

\subsection{Güç}
\[ y(t)=A_c\cos\big(2\pi f_c t+\phi(t)\big) \;\to\; \kutula{P_y=\frac{A_c^2}{2}} \qquad \text{güç hep buna eşit} \]

\subsection{BG}
\[ y(t)=\mathrm{Re}\{\hat y(t)\},\qquad \hat y(t)=A_c\, e^{j(2\pi f_c t+\phi(t))} \qquad [\,2\pi f_c t+\phi(t)=\theta(t)\,] \]
\[ \hat y(t)=A_c\, e^{j2\pi f_c t}\, e^{j\phi(t)} \]
\begin{align*}
e^{x}=\sum_{k=0}^{\infty}\frac{x^k}{k!} \;\to\; e^{j\phi(t)}&=1+j\phi(t)-\frac{\phi^2(t)}{2}-j\frac{\phi^3(t)}{3!}\cdots\\
&\quad+j^n\frac{\phi^n(t)}{n!}+\cdots
\end{align*}
\begin{align*}
\hat y(t)&=A_c\,[\cos(2\pi f_c t)+j\sin(2\pi f_c t)]\\
&\quad\cdot\Big[1+j\phi(t)-\frac{\phi^2(t)}{2}+\cdots+j^n\frac{\phi^n(t)}{n!}+\cdots\Big]
\end{align*}
\begin{align*}
y(t)=\mathrm{Re}\{\hat y(t)\}&=A_c\Big[\cos(2\pi f_c t)-\phi(t)\sin(2\pi f_c t)-\frac{\phi^2(t)}{2!}\cos(2\pi f_c t)\\
&\qquad+\frac{\phi^3(t)}{3!}\sin(2\pi f_c t)+\cdots\Big]
\end{align*}

\sayfa{6}
\[ \mathrm{BG}\{m(t)\}=B\ \mathrm{Hz} \;\Rightarrow\; \mathrm{BG}\{\phi(t)\}=B\ \mathrm{Hz} \]
\begin{align*}
&\hookrightarrow\ k_p\, m(t)\ \text{(PM)}\\
&\hookrightarrow\ 2\pi k_f\int m(\tau)\,d\tau\ \text{(FM)}
\end{align*}
\[ \mathrm{BG}\{\phi^2(t)\}=2B \;\longrightarrow\; \mathrm{BG}\{\phi^n(t)\}=nB \]
\begin{align*}
\mathrm{BG}\{y(t)\}&\to\infty\ \text{(Teorik)}\\
&\to\ \text{BG Teorikte sonsuzdur ancak güç zayıflar}
\end{align*}

\subsection{Açı modülasyonu Efektif BG}
Carson kuralı: $\mathrm{BG}=2B(\beta+1)$, $B=\mathrm{BG}\{m(t)\}$

FM(FRM):
\[ \beta=\frac{\Delta f}{B},\quad \Delta f=k_f\, m_p,\quad m_p=\max|m(t)| \]
($\beta$: sapma oranı (deviation Ratio); $\Delta f$: Taşıyıcıdan max sapma miktarı (Max carrier frequency deviation), $f_i(t)=f_c+k_f\,m(t)$)

PM(FZM):
\[ \beta=\frac{\Delta f}{B},\quad \Delta f=\frac{k_p}{2\pi}\,\frac{dm_p(t)}{dt},\quad \frac{dm_p(t)}{dt}=\max\left|\frac{dm(t)}{dt}\right| \]
\[ \Big(f_i(t)=f_c+\frac{k_p}{2\pi}\frac{dm(t)}{dt}\Big) \]

\begin{dikkat}
Not: kitapta buradaki formül yanlıştır. (O formül özel bir durumun formülü). Genelleştirilmiş için buradaki formülleri kullanacağız.
\end{dikkat}
\[ \theta(t)=2\pi f_c t+\overset{\text{faz}}{\phi(t)} \]
\[ \Delta\phi=\max|\phi(t)| \qquad \text{maks faz sapması (max phase deviation)} \]
FM(FRM):
\[ \Delta\phi=2\pi k_f\,\max\left|\int m(\tau)\,d\tau\right| \]
PM(FZM):
\[ \Delta\phi=k_p\, m_p,\qquad m_p=\max|m(t)| \]

\begin{ornek}
$m(t)=10\,\mathrm{sinc}(10^4 t)$ $\to$ Find the transmission BW of an FM modulated signal with $k_f=4000$
\end{ornek}
\[ m(t)=10\,\mathrm{sinc}(10^4 t) \;\to\; M(f)=\underset{A=\frac{1}{10^3}}{A}\,\Pi\!\left(\frac{f}{\tau}\right)=\frac{1}{10^3}\,\Pi\!\left(\frac{f}{10^4}\right) \]
\elcizimi[0.4]{sekiller/Ders11-s6-1.png}[6]
\[ B=5\ \mathrm{kHz}\qquad \beta=\frac{\Delta f}{B},\quad \Delta f=k_f\, m_p=4000\cdot10=40\ \mathrm{kHz} \]
\[ \beta=\frac{40\mathrm{k}}{5\mathrm{k}}=8 \;\Rightarrow\; \mathrm{BG}=2B(\beta+1)=2\cdot 5\mathrm{k}\,(8+1)=90\ \mathrm{kHz} \]
(90 kHz sonucundan okla yukarıdaki not:) mesajın BG'si 5k iken FRM'ninki 90k da çok çok büyük

\sayfa{7}
\begin{itemize}
\item Lineer değil (nonlinear)
\end{itemize}
\[ \phi(t)=\begin{cases} k_p\, m(t)\\ 2\pi k_f\displaystyle\int m(\tau)\,d\tau \end{cases} \]
\begin{align*}
A_c\cos\big(2\pi f_c t+a_1\phi_1(t)+a_2\phi_2(t)\big)&\neq A_c\cos\big(2\pi f_c t+a_1\phi_1(t)\big)\\
&\quad+A_c\cos\big(2\pi f_c t+a_2\phi_2(t)\big)
\end{align*}
\[ k_f\ll k_p\ll 1 \;\Rightarrow\; |\phi(t)|\ll1 \;\Rightarrow\; \phi^n(t)\approx0,\ n\ge2 \]
\[ y(t)\approx A_c\big[\cos(2\pi f_c t)-\phi(t)\sin(2\pi f_c t)\big] \]
$\hookrightarrow$ Dar bantlı açı modülasyonu (Narrow band angle mod)

BG az olan sabit zarflı bir modülasyon tekniği elde ettik

\subsection{Narrow Band angle modulation}
\[ \left.\begin{array}{l}k_f\ll1\\ k_p\ll1\end{array}\right\}\Rightarrow|\phi(t)|\ll1\Rightarrow\begin{cases}\cos(\phi(t))\approx1\\ \sin(\phi(t))\approx\phi(t)\end{cases} \]
\[ \cos(a+b)=\cos a\cdot\cos b-\sin a\cdot\sin b \]
\begin{align*}
y(t)&=A_c\cos\big(\underset{a}{2\pi f_c t}+\underset{b}{\phi(t)}\big)\\
&=A_c\big[\cos(2\pi f_c t)\cdot\underset{\approx1}{\cos(\phi(t))}-\sin(2\pi f_c t)\underset{\approx\phi(t)}{\sin\phi(t)}\big]
\end{align*}
\begin{align*}
y(t)&=\underbrace{A_c\big[\cos(2\pi f_c t)-\underset{\mathrm{BG}=B\,\mathrm{Hz}}{\phi(t)}\sin(2\pi f_c t)\big]}_{\mathrm{BG}=2B\ \mathrm{Hz}}\\
&\longrightarrow\ \phi(t)=k_p\,m(t)\ \text{(NBPM)}\\
&\longrightarrow\ \phi(t)=2\pi k_f\int m(\tau)\,d\tau\ \text{(NBFM)}
\end{align*}
\[ \Rightarrow\ \mathrm{BG}\{y(t)\}=2B\ \mathrm{Hz} \]

\subsection{NBPM Modulator}
\begin{center}
\begin{tikzpicture}[blok/.style={draw=murekkep,thick,align=center,text=murekkep,minimum height=9mm,inner sep=3pt}]
\node[left] (m) at (0,0) {$m(t)$};
\node[blok] (dsb) at (2.2,0) {DSB-SC-AM\\ Modulator};
\node[draw=murekkep,thick,circle,inner sep=1pt] (sum) at (7.2,0) {$+$};
\node[blok,minimum width=12mm] (ph) at (2.2,-2.3) {$\pi/2$};
\node[draw=murekkep,thick,circle,minimum size=9mm,inner sep=0pt] (osc) at (2.2,-4.3) {$\sim$};
\draw[->] (m.east) -- (dsb.west);
\draw[->] (dsb.east) -- (sum.west) node[midway,above] {\scriptsize $-A_ck_pm(t)\sin(2\pi f_ct)$};
\draw[->] (sum.east) -- ++(1.4,0);
\node[anchor=south west,inner sep=1pt] at ([xshift=1mm]sum.east) {\scriptsize $A_c\cos(2\pi f_ct)-A_ck_pm(t)\sin(2\pi f_ct)$};
\node[anchor=north west,inner sep=1pt] at ([xshift=8mm]sum.east) {\scriptsize $y_{\mathrm{NBPM}}(t)$};
\draw[->] (osc.north) -- (ph.south);
\draw[->] (ph.north) -- (dsb.south) node[midway,right] {\scriptsize $-A_c\sin(2\pi f_ct)$};
\draw[->] (osc.east) -| (sum.south);
\node[below] at (osc.south) {\scriptsize $A_c\cos(2\pi f_ct)$};
\end{tikzpicture}
\end{center}

\sayfa{8}
\subsection{NBFM Modulator}
\begin{center}
\begin{tikzpicture}[blok/.style={draw=murekkep,thick,align=center,text=murekkep,minimum height=9mm,inner sep=3pt}]
\node[left] (m) at (0,0) {$m(t)$};
\node[blok] (int) at (1.6,0) {$\displaystyle\int m(\tau)\,d\tau$};
\node[blok] (dsb) at (4.6,0) {DSB-SC-AM\\ Modulator};
\node[draw=murekkep,thick,circle,inner sep=1pt] (sum) at (9.9,0) {$+$};
\node[blok,minimum width=12mm] (ph) at (4.6,-2.3) {$\pi/2$};
\node[draw=murekkep,thick,circle,minimum size=9mm,inner sep=0pt] (osc) at (4.6,-4.3) {$\sim$};
\draw[->] (m.east) -- (int.west);
\draw[->] (int.east) -- (dsb.west);
\draw[->] (dsb.east) -- (sum.west) node[midway,above=1mm] {\scriptsize $-A_c\Big(\overbrace{2\pi k_f\int m(\tau)\,d\tau}^{\phi(t)}\Big)\sin(2\pi f_ct)$};
\draw[->] (sum.east) -- ++(1.0,0) node[right] {$y_{\mathrm{NBFM}}(t)$};
\draw[->] (osc.north) -- (ph.south);
\draw[->] (ph.north) -- (dsb.south) node[midway,right] {\scriptsize $-A_c\sin(2\pi f_ct)$};
\draw[->] (osc.east) -| (sum.south);
\node[below] at (osc.south) {\scriptsize $A_c\cos(2\pi f_ct)$};
\end{tikzpicture}
\end{center}

Carson kuralı $\mathrm{BG}=2B(\beta+1)$
\[
\left.\begin{aligned}
\text{FM: }\Delta f&=k_f\,m_p \xrightarrow{\ k_f\ll1\ } \Delta f\ll\\
\text{PM: }\Delta f&=\frac{k_p}{2\pi}\frac{dm_p}{dt} \xrightarrow{\ k_p\ll1\ } \Delta f\ll
\end{aligned}\right\}
\begin{aligned}
\beta&=\frac{\Delta f}{B}\ll1\\
\mathrm{BG}&=2B(\beta+1)\approx2B
\end{aligned}
\]
\[ \text{(NB }\overset{\text{Angle}}{\text{A}}\text{M)} \]

\subsection{Wide Band Angle modulation (WBAM)}
\[
\left.\begin{aligned}
\text{FM: }\Delta f&=k_f\,m_p\\
\text{PM: }\Delta f&=\frac{k_p}{2\pi}\frac{dm_p}{dt}
\end{aligned}\right\}
\begin{array}{l}k_f\gg\\ k_p\gg\end{array}
\;\Rightarrow\;\Delta f\gg\;\Rightarrow\;\beta=\frac{\Delta f}{B}\gg1
\]
\[ \mathrm{BG}=2B(\beta+1)\approx2B\beta\approx2\Delta f \]

\end{document}
